% =============================================================================
% architekturprinzip.tex – Robuste Embedded-Software
% Einbindung: % =============================================================================
% architekturprinzip.tex – Robuste Embedded-Software
% Einbindung: % =============================================================================
% architekturprinzip.tex – Robuste Embedded-Software
% Einbindung: % =============================================================================
% architekturprinzip.tex – Robuste Embedded-Software
% Einbindung: \input{content/architekturprinzip}
% =============================================================================

\section{Architekturprinzipien -- Robuste Embedded-Software}
\label{sec:architekturprinzipien}

% -----------------------------------------------------------------------------
\subsection{Begriffsrahmen und Zielsetzung}
\label{subsec:arch-begriffsrahmen}

Die Selection-Panel-Firmware steuert 10--100 Buttons und LEDs über Schieberegister
(CD4021, 74HC595). Die Architektur muss drei Anforderungen gleichzeitig erfüllen:
Korrektheit (One-Hot-Invariante), Wartbarkeit (klare Zuständigkeiten), Skalierbarkeit
(10 $\to$ 100 ohne Code-Änderungen in der Logik).

\begin{table}[H]
\centering
\caption{Zentrale Begriffe der Firmware-Architektur}
\label{tab:arch-begriffe}
\begin{tabularx}{\textwidth}{@{}l X@{}}
\toprule
\textbf{Begriff} & \textbf{Definition} \\
\midrule
SSOT       & Single Source of Truth -- eine Konfiguration existiert an genau einer Stelle \\
Invariante & Bedingung, die zu jedem Zeitpunkt gelten muss \\
One-Hot    & Genau ein Bit gesetzt (oder keins), niemals mehrere \\
Pipeline   & Sequenz von Verarbeitungsstufen mit definierten Ein-/Ausgängen \\
Edge       & Zustandsübergang (Flanke), nicht der Zustand selbst \\
HAL        & Hardware Abstraction Layer -- trennt Hardware-Zugriff von Logik \\
\bottomrule
\end{tabularx}
\end{table}

\begin{tipbox}[Ziel]
Stabiler Datenfluss, testbare Logik, keine Seiteneffekte.
\end{tipbox}

% -----------------------------------------------------------------------------
\subsection{SSOT -- Hardware-Konfiguration ist kein \glqq Detail\grqq}
\label{subsec:ssot-prinzip}

Wenn Pins, Zählerstände oder SPI-Parameter an mehreren Stellen definiert sind,
entstehen Inkonsistenzen. Du änderst einen Pin in \code{main.cpp}, vergisst aber
den Treiber -- das System verhält sich unerklärlich.

\textbf{Regel:} Jede Konfiguration existiert an genau einer Stelle. Alle anderen
Dateien referenzieren diese Quelle.

\begin{lstlisting}[style=shell,numbers=none,caption={SSOT-Verzeichnisstruktur},label={lst:ssot-struktur}]
firmware/include/hw/
|-- board_pins.h    # Pin-Definitionen (SSOT)
|-- io_counts.h     # Counts, Masks (SSOT)
`-- spi_params.h    # SPI-Parameter (SSOT)
\end{lstlisting}

\paragraph{Beispiel aus io\_counts.h}

\begin{lstlisting}[style=arduino,caption={SSOT für Counts und Masken},label={lst:io-counts}]
#define BTN_COUNT       10
#define LED_COUNT       10

#define BTN_BYTES       ((BTN_COUNT + 7) / 8)  // = 2
#define LED_BYTES       ((LED_COUNT + 7) / 8)  // = 2

// Maske fuer ungueltige Bits im letzten Byte
#define LED_LAST_MASK   0xC0  // Bei 10 LEDs: nur Bit 7+6 gueltig
\end{lstlisting}

\begin{infobox}[Konsequenz]
Wenn du von 10 auf 100 Kanäle skalierst, änderst du \code{BTN\_COUNT} und
\code{LED\_COUNT} an dieser einen Stelle. Die Makros \code{BTN\_BYTES} und
\code{LED\_LAST\_MASK} berechnen sich automatisch neu.
\end{infobox}

% -----------------------------------------------------------------------------
\subsection{Schichtenmodell -- Abhängigkeiten zeigen nach unten}
\label{subsec:schichtenmodell}

Ohne klare Schichtung entsteht \glqq Spaghetti-Code\grqq: Der Treiber ruft Logik auf,
die Logik greift auf Hardware zu, die App dupliziert Regeln. Änderungen haben
unvorhersehbare Auswirkungen.

\textbf{Regel:} Abhängigkeiten sind gerichtet. Höhere Schichten dürfen niedrigere
nutzen, aber nicht umgekehrt.

\begin{figure}[H]
\centering
\begin{tikzpicture}[
  layer/.style={draw, rectangle, minimum width=11cm, minimum height=0.9cm,
                font=\ttfamily\small, fill=white},
  label/.style={font=\sffamily\footnotesize, text=MutedText}
]
  \node[layer, fill=ArduinoTeal!15] (app) at (0,4) {app/ \hfill Orchestrierung (Loop/Task, Scheduling)};
  \node[layer, fill=ArduinoTeal!10] (logic) at (0,3) {logic/ \hfill Regeln (Debounce, Selection, Pipeline)};
  \node[layer, fill=SuccessGreen!10] (drivers) at (0,2) {drivers/ \hfill IC-spezifisch (CD4021 read, HC595 write)};
  \node[layer, fill=SuccessGreen!15] (hal) at (0,1) {hal/ \hfill SPI-Bus Zugriff (Transfer/Settings)};
  \node[layer, fill=MutedText!10] (hw) at (0,0) {include/hw/* \hfill SSOT (Pins/Counts/SPI-Params)};

  \draw[-{Stealth}, thick] (app.south west) ++(0.3,0) -- ++(0,-0.15);
  \draw[-{Stealth}, thick] (logic.south west) ++(0.3,0) -- ++(0,-0.15);
  \draw[-{Stealth}, thick] (drivers.south west) ++(0.3,0) -- ++(0,-0.15);
  \draw[-{Stealth}, thick] (hal.south west) ++(0.3,0) -- ++(0,-0.15);
\end{tikzpicture}
\caption{Schichtenmodell der Firmware}
\label{fig:schichtenmodell}
\end{figure}

\paragraph{Verzeichnisstruktur}

\begin{lstlisting}[style=shell,numbers=none,caption={Firmware-Verzeichnisstruktur},label={lst:firmware-struktur}]
firmware/src/
|-- hal/
|   `-- spi_bus.cpp         # SPI-Initialisierung
|-- drivers/
|   |-- cd4021.cpp          # Button-Input (liest Bytes)
|   `-- hc595.cpp           # LED-Output (schreibt Bytes)
|-- logic/
|   |-- debounce.cpp        # Entprellung -> Press-Edges
|   |-- selection.cpp       # Selection-Policy -> selectedId
|   `-- io_pipeline.cpp     # Orchestriert Logik-Module
`-- app/
    |-- io_loop.cpp         # Arduino-Variante
    `-- io_task.cpp         # FreeRTOS-Variante
\end{lstlisting}

\begin{table}[H]
\centering
\caption{Schicht-Verantwortlichkeiten}
\label{tab:schicht-regeln}
\begin{tabularx}{\textwidth}{@{}l X X@{}}
\toprule
\textbf{Schicht} & \textbf{Darf} & \textbf{Darf nicht} \\
\midrule
\code{drivers/} & Bytes lesen/schreiben        & Debounce-Entscheidungen treffen \\
\code{logic/}   & IDs und Zustände verarbeiten & Pins kennen, SPI aufrufen \\
\code{app/}     & Schichten orchestrieren      & Regeln duplizieren \\
\bottomrule
\end{tabularx}
\end{table}

\begin{tipbox}[Warum ist das wichtig?]
Die \code{logic/}-Schicht ist hardware-unabhängig. Du kannst sie auf dem Host-PC
unit-testen, ohne ESP32, ohne SPI, ohne echte Hardware.
\end{tipbox}

% -----------------------------------------------------------------------------
\subsection{Datenfluss ist eine Pipeline}
\label{subsec:pipeline}

Wenn IO-Code \glqq verstreut\grqq\ ist -- hier ein LED-Update, dort ein Button-Check --
wird das System unvorhersehbar. Die One-Hot-Invariante kann brechen, weil zwei
Codepfade unabhängig LEDs setzen.

\textbf{Regel:} Der gesamte IO-Zyklus folgt einer klaren Sequenz:

\begin{lstlisting}[style=base,numbers=none]
read(raw) -> normalize/logic(update) -> derive outputs -> write(outputs)
\end{lstlisting}

\begin{figure}[H]
\centering
\begin{tikzpicture}[
  box/.style={draw, rectangle, rounded corners=2pt, minimum height=0.8cm,
              font=\ttfamily\small, fill=white},
  pipeline/.style={draw, rectangle, rounded corners=3pt, minimum width=10cm,
                   minimum height=1.8cm, fill=ArduinoTeal!8},
  arrow/.style={-{Stealth[length=2mm]}, thick}
]
  % CD4021
  \node[box, fill=SuccessGreen!15] (cd4021) at (0,3) {CD4021::read(rawButtons)};

  % Pipeline container
  \node[pipeline] (pipe) at (0,1.2) {};
  \node[font=\sffamily\bfseries\small, anchor=north] at (0,2) {io\_pipeline};

  % Pipeline stages
  \node[box] (debounce) at (-3.5,1.2) {Debounce\\(Edges)};
  \node[box] (guard) at (0,1.2) {Press-Guard\\(optional)};
  \node[box] (selection) at (3.5,1.2) {Selection\\(Policy)};

  \draw[arrow] (debounce) -- (guard);
  \draw[arrow] (guard) -- (selection);

  % selectedId
  \node[font=\ttfamily\small] (selid) at (0,-0.3) {selectedId};

  % setOneHot
  \node[box, fill=ArduinoTeal!15] (onehot) at (0,-1.3) {setOneHot(selectedId) $\to$ ledState[]};

  % HC595
  \node[box, fill=SuccessGreen!15] (hc595) at (0,-2.5) {HC595::write(ledState)};

  % Arrows
  \draw[arrow] (cd4021) -- (pipe.north);
  \draw[arrow] (pipe.south) -- (selid);
  \draw[arrow] (selid) -- (onehot);
  \draw[arrow] (onehot) -- (hc595);
\end{tikzpicture}
\caption{Datenfluss-Pipeline im Selection Panel}
\label{fig:pipeline}
\end{figure}

\paragraph{Was passiert in jeder Stufe?}

\begin{enumerate}
  \item \textbf{CD4021::read()} -- Liest rohe Bytes vom Schieberegister (active-low, 0~=~gedrückt)
  \item \textbf{Debounce} -- Filtert Prellen, erzeugt stabile Press-Edges
  \item \textbf{Press-Guard} -- Verhindert versehentliche Doppel-Presses (optional)
  \item \textbf{Selection} -- Wendet Policy an: \glqq Welche LED soll leuchten?\grqq
  \item \textbf{setOneHot()} -- Wandelt \code{selectedId} in Byte-Array um
  \item \textbf{HC595::write()} -- Schiebt Bytes zum Schieberegister
\end{enumerate}

\begin{warnbox}[Konsequenz]
LED-Zustände entstehen an genau einer Stelle (\code{setOneHot}). Kein anderer Code
darf LEDs direkt setzen -- das würde die Pipeline umgehen und könnte One-Hot brechen.
\end{warnbox}

% -----------------------------------------------------------------------------
\subsection{One-Hot ist eine systemweite Invariante}
\label{subsec:one-hot}

Bei einem Panel mit 100~LEDs wäre es fatal, wenn versehentlich zwei LEDs gleichzeitig
leuchten oder ein ungültiger Index adressiert wird.

\textbf{Regel:} One-Hot ist ein Vertrag:
\begin{itemize}
  \item \code{selectedId == 0} $\to$ alle LEDs aus
  \item \code{selectedId} $\in [1..N]$ $\to$ genau eine LED an
  \item Andere Zustände sind verboten
\end{itemize}

\paragraph{Implementierung der Invariante}

\begin{lstlisting}[style=arduino,caption={One-Hot Implementierung},label={lst:one-hot}]
void setOneHot(uint8_t selectedId, uint8_t* ledState, size_t ledBytes) {
    // 1. Alle LEDs aus (clear)
    memset(ledState, 0x00, ledBytes);

    // 2. Genau eine LED setzen (wenn selectedId > 0)
    if (selectedId >= 1 && selectedId <= LED_COUNT) {
        uint8_t byteIdx = (selectedId - 1) / 8;
        uint8_t bitIdx  = 7 - ((selectedId - 1) % 8);  // MSB-first
        ledState[byteIdx] |= (1 << bitIdx);
    }

    // 3. Ungueltige Bits im letzten Byte maskieren
    ledState[ledBytes - 1] &= LED_LAST_MASK;
}
\end{lstlisting}

\begin{infobox}[Warum Schritt 3?]
Bei 10~LEDs brauchst du 2~Bytes (16~Bits), aber nur 10 sind gültig. Die Bits 5--0
im zweiten Byte könnten durch Bugs oder Speicherfehler gesetzt sein. Die Maske
garantiert, dass sie niemals an den HC595 gelangen.
\end{infobox}

\paragraph{Beispiel (10 LEDs, MSB-first)}

\begin{table}[H]
\centering
\caption{One-Hot Bit-Mapping für 10 LEDs}
\label{tab:one-hot-mapping}
\begin{tabular}{@{}c l l l@{}}
\toprule
\textbf{selectedId} & \textbf{Byte[0]} & \textbf{Byte[1]} & \textbf{Hex} \\
\midrule
0 (aus) & \code{00000000} & \code{00000000} & \code{0x00 0x00} \\
1       & \code{10000000} & \code{00000000} & \code{0x80 0x00} \\
8       & \code{00000001} & \code{00000000} & \code{0x01 0x00} \\
9       & \code{00000000} & \code{10000000} & \code{0x00 0x80} \\
10      & \code{00000000} & \code{01000000} & \code{0x00 0x40} \\
\bottomrule
\end{tabular}
\end{table}

% -----------------------------------------------------------------------------
\subsection{Ereignisse statt Pegel -- Buttons sind Edge-getrieben}
\label{subsec:edge-driven}

Wenn die Logik auf \glqq Button ist gedrückt\grqq\ (Pegel) reagiert, feuert sie bei
jedem Scan-Zyklus. Bei \SI{10}{\milli\second} Scan-Intervall wären das 100~Events
pro Sekunde, solange du die Taste hältst.

\textbf{Regel:} UI-Aktionen basieren auf Press-Edges (Flanken), nicht auf Pegeln.

\begin{lstlisting}[style=base,numbers=none]
CD4021 (Pegel)      Debounce            Selection
     |                  |                   |
  active-low    ->   pressed edges    ->   reagiert nur
  (0=gedrueckt)     (stabil+Flanke)       auf Edges
\end{lstlisting}

\paragraph{Debounce-Algorithmus (vereinfacht)}

\begin{lstlisting}[style=arduino,caption={Debounce-Implementierung},label={lst:debounce}]
typedef struct {
    uint8_t stable;      // Letzter stabiler Zustand
    uint8_t counter;     // Zaehler fuer Entprellung
} DebounceState;

bool debounce_update(DebounceState* state, uint8_t raw, uint8_t threshold) {
    if (raw == state->stable) {
        state->counter = 0;
        return false;  // Keine Aenderung
    }

    state->counter++;
    if (state->counter >= threshold) {
        bool wasPressed = (state->stable != 0);
        bool isPressed  = (raw != 0);
        state->stable = raw;
        state->counter = 0;

        // Press-Edge: war nicht gedrueckt, ist jetzt gedrueckt
        return (!wasPressed && isPressed);
    }
    return false;
}
\end{lstlisting}

\begin{tipbox}[Konsequenz]
Die Selection-Logik sieht nur \glqq Button 7 wurde gerade gedrückt\grqq, nicht
\glqq Button 7 ist seit \SI{500}{\milli\second} gedrückt\grqq. Das verhindert
Event-Spam und macht das Verhalten vorhersagbar.
\end{tipbox}

% -----------------------------------------------------------------------------
\subsection{Determinismus -- Timing bleibt stabil}
\label{subsec:determinismus}

Im Embedded-Kontext kann ein einzelner \code{Serial.printf()} im Hot-Path das
gesamte Timing zerstören. UART bei \SI{115200}{Baud} braucht etwa \SI{87}{\micro\second}
pro Zeichen -- bei 50~Zeichen sind das \SI{4,3}{\milli\second}, fast die Hälfte
eines \SI{10}{\milli\second}-Scan-Intervalls.

\textbf{Regel:} Der Scan-Pfad läuft ohne Blockierungen. Logging erfolgt nur bei
Zustandsänderungen.

\begin{table}[H]
\centering
\caption{Erlaubte und verbotene Operationen im Hot-Path}
\label{tab:timing-regeln}
\begin{tabularx}{\textwidth}{@{}X X@{}}
\toprule
\textbf{Erlaubt} & \textbf{Verboten} \\
\midrule
Logging bei Selection-Change      & Logging pro Scan-Zyklus \\
\code{vTaskDelayUntil()} (FreeRTOS) & Busy-loops mit \code{delay()} \\
Kurze SPI-Transfers ($\sim$\SI{10}{\micro\second}) & Lange Wartezeiten im Hot-Path \\
\bottomrule
\end{tabularx}
\end{table}

\paragraph{Timing-Budget (10 Kanäle)}

\begin{table}[H]
\centering
\caption{Timing-Budget pro Scan-Zyklus}
\label{tab:timing-budget}
\begin{tabular}{@{}l r@{}}
\toprule
\textbf{Operation} & \textbf{Typische Dauer} \\
\midrule
CD4021 read (2~Bytes)              & $\sim$\SI{20}{\micro\second} \\
Debounce + Selection               & $\sim$\SI{5}{\micro\second} \\
HC595 write (2~Bytes)              & $\sim$\SI{20}{\micro\second} \\
\midrule
\textbf{Gesamt pro Zyklus}         & $\sim$\SI{50}{\micro\second} \\
\textbf{Verfügbar bei \SI{10}{\milli\second} Intervall} & \SI{10000}{\micro\second} \\
\bottomrule
\end{tabular}
\end{table}

\begin{infobox}[Konsequenz]
Bei 10~Kanälen hast du reichlich Reserve. Bei 100~Kanälen (13~Bytes,
$\sim$\SI{130}{\micro\second} Transfer) bleibt das Timing stabil, solange du
keine blockierenden Operationen einbaust.
\end{infobox}

% -----------------------------------------------------------------------------
\subsection{Skalierbarkeit -- Byte-orientiert und maskenbasiert}
\label{subsec:skalierbarkeit}

Code mit hartkodierten Werten wie \code{uint8\_t buttons[2]} oder Schleifen
\code{for (int i = 0; i < 10; i++)} muss bei Skalierung manuell angepasst werden.

\textbf{Regel:} Algorithmen arbeiten mit \code{*\_COUNT} und \code{*\_BYTES},
nicht mit Literalen.

\begin{lstlisting}[style=arduino,caption={Skalierbare vs.\ hartkodierte Implementierung},label={lst:skalierbar}]
// Schlecht: hartkodiert
uint8_t ledState[2];
for (int i = 0; i < 10; i++) { ... }

// Gut: skalierbar
uint8_t ledState[LED_BYTES];
for (int i = 0; i < LED_COUNT; i++) { ... }
\end{lstlisting}

\paragraph{Automatische Berechnung}

\begin{lstlisting}[style=arduino,caption={Automatische Maskenberechnung für 100 Kanäle},label={lst:auto-mask}]
#define BTN_COUNT   100
#define BTN_BYTES   ((BTN_COUNT + 7) / 8)  // = 13

#define LED_COUNT   100
#define LED_BYTES   ((LED_COUNT + 7) / 8)  // = 13

// Maske: Bei 100 LEDs sind im letzten Byte nur Bits 7-4 gueltig (4 LEDs)
// 100 % 8 = 4 -> Bits 7,6,5,4 -> Maske = 0xF0
#define LED_LAST_MASK  (0xFF << (8 - (LED_COUNT % 8))) & 0xFF
\end{lstlisting}

\begin{tipbox}[Konsequenz]
Die Änderung \code{BTN\_COUNT=10} $\to$ \code{BTN\_COUNT=100} erfordert keine
Anpassung in der Logik. Arrays, Schleifen und Masken passen sich automatisch an.
\end{tipbox}

% -----------------------------------------------------------------------------
\subsection{Testbarkeit ist ein Architekturziel}
\label{subsec:testbarkeit}

Wenn Logik und Hardware verwoben sind, kannst du nur am echten Gerät testen.
Jeder Test braucht Flashen, Verkabeln, manuelles Drücken -- langsam und fehleranfällig.

\textbf{Regel:} Logik muss isoliert testbar sein, Hardware-Pfad end-to-end.

\begin{table}[H]
\centering
\caption{Test-Environments in platformio.ini}
\label{tab:test-environments}
\begin{tabularx}{\textwidth}{@{}l l X@{}}
\toprule
\textbf{Environment} & \textbf{Typ} & \textbf{Was wird getestet} \\
\midrule
\code{seeed\_xiao\_esp32s3\_test} & Unit & Nur \code{logic/} (Debounce, Selection, Pipeline) \\
\code{seeed\_xiao\_esp32s3\_e2e}  & E2E  & \code{logic/} + \code{drivers/} + \code{hal/} mit echter Hardware \\
\bottomrule
\end{tabularx}
\end{table}

\paragraph{Unit-Test-Beispiel (Debounce)}

\begin{lstlisting}[style=arduino,caption={Unit-Test für Debounce},label={lst:debounce-test}]
void test_debounce_filters_noise() {
    DebounceState state = {0};

    // Simuliere Prellen: 0 -> 1 -> 0 -> 1 -> 1 -> 1
    TEST_ASSERT_FALSE(debounce_update(&state, 1, 3));  // Zaehler = 1
    TEST_ASSERT_FALSE(debounce_update(&state, 0, 3));  // Reset
    TEST_ASSERT_FALSE(debounce_update(&state, 1, 3));  // Zaehler = 1
    TEST_ASSERT_FALSE(debounce_update(&state, 1, 3));  // Zaehler = 2
    TEST_ASSERT_TRUE(debounce_update(&state, 1, 3));   // Zaehler = 3 -> Edge!
}
\end{lstlisting}

\begin{infobox}[Konsequenz]
Unit-Tests laufen in Sekunden, nicht Minuten. Du kannst Debounce-Schwellenwerte
experimentell optimieren, ohne jedes Mal zu flashen.
\end{infobox}

% -----------------------------------------------------------------------------
\subsection{Änderungen haben Pflichtfolgen}
\label{subsec:pflichtfolgen}

\textbf{Regel:} Bestimmte Änderungen erzwingen bestimmte Checks.

\begin{table}[H]
\centering
\caption{Pflicht-Checks nach Änderungstyp}
\label{tab:pflicht-checks}
\begin{tabularx}{\textwidth}{@{}l X@{}}
\toprule
\textbf{Änderung} & \textbf{Pflicht} \\
\midrule
Pins/SPI/Mapping            & \code{wiring.md} aktualisieren + E2E-Test \\
Debounce/Selection/Pipeline & Unit-Tests (\code{seeed\_xiao\_esp32s3\_test}) \\
Treiber/HAL                 & E2E-Test + Timing-Messung \\
Counts (10 $\to$ 100)       & Unit-Tests + E2E-Test + Timing-Prüfung \\
\bottomrule
\end{tabularx}
\end{table}

\paragraph{Testbefehle}

\begin{lstlisting}[style=shell,caption={Testbefehle für Unit- und E2E-Tests},label={lst:testbefehle}]
# Unit-Tests (Logik)
pio test -e seeed_xiao_esp32s3_test -v

# E2E-Hardware-Test
pio test -e seeed_xiao_esp32s3_e2e -f test_e2e_hw -v
\end{lstlisting}

% -----------------------------------------------------------------------------
\subsection{Minimal-Checkliste}
\label{subsec:checkliste}

Bevor du einen Change committest:

\begin{itemize}
  \item[$\square$] SSOT nur in \code{include/hw/*}
  \item[$\square$] Datenfluss bleibt: read $\to$ update $\to$ write
  \item[$\square$] One-Hot entsteht in genau einem Pfad (clear $\to$ set)
  \item[$\square$] Letztes Byte maskiert (\code{*\_LAST\_MASK})
  \item[$\square$] \code{logic/} ist hardware-unabhängig
  \item[$\square$] Unit-Tests grün (\code{seeed\_xiao\_esp32s3\_test})
  \item[$\square$] Bei HW-relevanten Changes: E2E grün (\code{test\_e2e\_hw})
\end{itemize}

% -----------------------------------------------------------------------------
\subsection{Siehe auch}
\label{subsec:arch-siehe-auch}

Weiterführende Dokumentation:

\begin{itemize}
  \item \code{firmware/docs/architecture.md} -- Schichtenmodell-Details
  \item \code{firmware/docs/overview.md} -- Bit-Mapping (MSB-first)
  \item \code{firmware/docs/wiring.md} -- Hardware-Verdrahtung
  \item \code{apps/bridge/docs/PROTOCOL.md} -- Serial/WebSocket-Protokoll
\end{itemize}

% =============================================================================

\section{Architekturprinzipien -- Robuste Embedded-Software}
\label{sec:architekturprinzipien}

% -----------------------------------------------------------------------------
\subsection{Begriffsrahmen und Zielsetzung}
\label{subsec:arch-begriffsrahmen}

Die Selection-Panel-Firmware steuert 10--100 Buttons und LEDs über Schieberegister
(CD4021, 74HC595). Die Architektur muss drei Anforderungen gleichzeitig erfüllen:
Korrektheit (One-Hot-Invariante), Wartbarkeit (klare Zuständigkeiten), Skalierbarkeit
(10 $\to$ 100 ohne Code-Änderungen in der Logik).

\begin{table}[H]
\centering
\caption{Zentrale Begriffe der Firmware-Architektur}
\label{tab:arch-begriffe}
\begin{tabularx}{\textwidth}{@{}l X@{}}
\toprule
\textbf{Begriff} & \textbf{Definition} \\
\midrule
SSOT       & Single Source of Truth -- eine Konfiguration existiert an genau einer Stelle \\
Invariante & Bedingung, die zu jedem Zeitpunkt gelten muss \\
One-Hot    & Genau ein Bit gesetzt (oder keins), niemals mehrere \\
Pipeline   & Sequenz von Verarbeitungsstufen mit definierten Ein-/Ausgängen \\
Edge       & Zustandsübergang (Flanke), nicht der Zustand selbst \\
HAL        & Hardware Abstraction Layer -- trennt Hardware-Zugriff von Logik \\
\bottomrule
\end{tabularx}
\end{table}

\begin{tipbox}[Ziel]
Stabiler Datenfluss, testbare Logik, keine Seiteneffekte.
\end{tipbox}

% -----------------------------------------------------------------------------
\subsection{SSOT -- Hardware-Konfiguration ist kein \glqq Detail\grqq}
\label{subsec:ssot-prinzip}

Wenn Pins, Zählerstände oder SPI-Parameter an mehreren Stellen definiert sind,
entstehen Inkonsistenzen. Du änderst einen Pin in \code{main.cpp}, vergisst aber
den Treiber -- das System verhält sich unerklärlich.

\textbf{Regel:} Jede Konfiguration existiert an genau einer Stelle. Alle anderen
Dateien referenzieren diese Quelle.

\begin{lstlisting}[style=shell,numbers=none,caption={SSOT-Verzeichnisstruktur},label={lst:ssot-struktur}]
firmware/include/hw/
|-- board_pins.h    # Pin-Definitionen (SSOT)
|-- io_counts.h     # Counts, Masks (SSOT)
`-- spi_params.h    # SPI-Parameter (SSOT)
\end{lstlisting}

\paragraph{Beispiel aus io\_counts.h}

\begin{lstlisting}[style=arduino,caption={SSOT für Counts und Masken},label={lst:io-counts}]
#define BTN_COUNT       10
#define LED_COUNT       10

#define BTN_BYTES       ((BTN_COUNT + 7) / 8)  // = 2
#define LED_BYTES       ((LED_COUNT + 7) / 8)  // = 2

// Maske fuer ungueltige Bits im letzten Byte
#define LED_LAST_MASK   0xC0  // Bei 10 LEDs: nur Bit 7+6 gueltig
\end{lstlisting}

\begin{infobox}[Konsequenz]
Wenn du von 10 auf 100 Kanäle skalierst, änderst du \code{BTN\_COUNT} und
\code{LED\_COUNT} an dieser einen Stelle. Die Makros \code{BTN\_BYTES} und
\code{LED\_LAST\_MASK} berechnen sich automatisch neu.
\end{infobox}

% -----------------------------------------------------------------------------
\subsection{Schichtenmodell -- Abhängigkeiten zeigen nach unten}
\label{subsec:schichtenmodell}

Ohne klare Schichtung entsteht \glqq Spaghetti-Code\grqq: Der Treiber ruft Logik auf,
die Logik greift auf Hardware zu, die App dupliziert Regeln. Änderungen haben
unvorhersehbare Auswirkungen.

\textbf{Regel:} Abhängigkeiten sind gerichtet. Höhere Schichten dürfen niedrigere
nutzen, aber nicht umgekehrt.

\begin{figure}[H]
\centering
\begin{tikzpicture}[
  layer/.style={draw, rectangle, minimum width=11cm, minimum height=0.9cm,
                font=\ttfamily\small, fill=white},
  label/.style={font=\sffamily\footnotesize, text=MutedText}
]
  \node[layer, fill=ArduinoTeal!15] (app) at (0,4) {app/ \hfill Orchestrierung (Loop/Task, Scheduling)};
  \node[layer, fill=ArduinoTeal!10] (logic) at (0,3) {logic/ \hfill Regeln (Debounce, Selection, Pipeline)};
  \node[layer, fill=SuccessGreen!10] (drivers) at (0,2) {drivers/ \hfill IC-spezifisch (CD4021 read, HC595 write)};
  \node[layer, fill=SuccessGreen!15] (hal) at (0,1) {hal/ \hfill SPI-Bus Zugriff (Transfer/Settings)};
  \node[layer, fill=MutedText!10] (hw) at (0,0) {include/hw/* \hfill SSOT (Pins/Counts/SPI-Params)};

  \draw[-{Stealth}, thick] (app.south west) ++(0.3,0) -- ++(0,-0.15);
  \draw[-{Stealth}, thick] (logic.south west) ++(0.3,0) -- ++(0,-0.15);
  \draw[-{Stealth}, thick] (drivers.south west) ++(0.3,0) -- ++(0,-0.15);
  \draw[-{Stealth}, thick] (hal.south west) ++(0.3,0) -- ++(0,-0.15);
\end{tikzpicture}
\caption{Schichtenmodell der Firmware}
\label{fig:schichtenmodell}
\end{figure}

\paragraph{Verzeichnisstruktur}

\begin{lstlisting}[style=shell,numbers=none,caption={Firmware-Verzeichnisstruktur},label={lst:firmware-struktur}]
firmware/src/
|-- hal/
|   `-- spi_bus.cpp         # SPI-Initialisierung
|-- drivers/
|   |-- cd4021.cpp          # Button-Input (liest Bytes)
|   `-- hc595.cpp           # LED-Output (schreibt Bytes)
|-- logic/
|   |-- debounce.cpp        # Entprellung -> Press-Edges
|   |-- selection.cpp       # Selection-Policy -> selectedId
|   `-- io_pipeline.cpp     # Orchestriert Logik-Module
`-- app/
    |-- io_loop.cpp         # Arduino-Variante
    `-- io_task.cpp         # FreeRTOS-Variante
\end{lstlisting}

\begin{table}[H]
\centering
\caption{Schicht-Verantwortlichkeiten}
\label{tab:schicht-regeln}
\begin{tabularx}{\textwidth}{@{}l X X@{}}
\toprule
\textbf{Schicht} & \textbf{Darf} & \textbf{Darf nicht} \\
\midrule
\code{drivers/} & Bytes lesen/schreiben        & Debounce-Entscheidungen treffen \\
\code{logic/}   & IDs und Zustände verarbeiten & Pins kennen, SPI aufrufen \\
\code{app/}     & Schichten orchestrieren      & Regeln duplizieren \\
\bottomrule
\end{tabularx}
\end{table}

\begin{tipbox}[Warum ist das wichtig?]
Die \code{logic/}-Schicht ist hardware-unabhängig. Du kannst sie auf dem Host-PC
unit-testen, ohne ESP32, ohne SPI, ohne echte Hardware.
\end{tipbox}

% -----------------------------------------------------------------------------
\subsection{Datenfluss ist eine Pipeline}
\label{subsec:pipeline}

Wenn IO-Code \glqq verstreut\grqq\ ist -- hier ein LED-Update, dort ein Button-Check --
wird das System unvorhersehbar. Die One-Hot-Invariante kann brechen, weil zwei
Codepfade unabhängig LEDs setzen.

\textbf{Regel:} Der gesamte IO-Zyklus folgt einer klaren Sequenz:

\begin{lstlisting}[style=base,numbers=none]
read(raw) -> normalize/logic(update) -> derive outputs -> write(outputs)
\end{lstlisting}

\begin{figure}[H]
\centering
\begin{tikzpicture}[
  box/.style={draw, rectangle, rounded corners=2pt, minimum height=0.8cm,
              font=\ttfamily\small, fill=white},
  pipeline/.style={draw, rectangle, rounded corners=3pt, minimum width=10cm,
                   minimum height=1.8cm, fill=ArduinoTeal!8},
  arrow/.style={-{Stealth[length=2mm]}, thick}
]
  % CD4021
  \node[box, fill=SuccessGreen!15] (cd4021) at (0,3) {CD4021::read(rawButtons)};

  % Pipeline container
  \node[pipeline] (pipe) at (0,1.2) {};
  \node[font=\sffamily\bfseries\small, anchor=north] at (0,2) {io\_pipeline};

  % Pipeline stages
  \node[box] (debounce) at (-3.5,1.2) {Debounce\\(Edges)};
  \node[box] (guard) at (0,1.2) {Press-Guard\\(optional)};
  \node[box] (selection) at (3.5,1.2) {Selection\\(Policy)};

  \draw[arrow] (debounce) -- (guard);
  \draw[arrow] (guard) -- (selection);

  % selectedId
  \node[font=\ttfamily\small] (selid) at (0,-0.3) {selectedId};

  % setOneHot
  \node[box, fill=ArduinoTeal!15] (onehot) at (0,-1.3) {setOneHot(selectedId) $\to$ ledState[]};

  % HC595
  \node[box, fill=SuccessGreen!15] (hc595) at (0,-2.5) {HC595::write(ledState)};

  % Arrows
  \draw[arrow] (cd4021) -- (pipe.north);
  \draw[arrow] (pipe.south) -- (selid);
  \draw[arrow] (selid) -- (onehot);
  \draw[arrow] (onehot) -- (hc595);
\end{tikzpicture}
\caption{Datenfluss-Pipeline im Selection Panel}
\label{fig:pipeline}
\end{figure}

\paragraph{Was passiert in jeder Stufe?}

\begin{enumerate}
  \item \textbf{CD4021::read()} -- Liest rohe Bytes vom Schieberegister (active-low, 0~=~gedrückt)
  \item \textbf{Debounce} -- Filtert Prellen, erzeugt stabile Press-Edges
  \item \textbf{Press-Guard} -- Verhindert versehentliche Doppel-Presses (optional)
  \item \textbf{Selection} -- Wendet Policy an: \glqq Welche LED soll leuchten?\grqq
  \item \textbf{setOneHot()} -- Wandelt \code{selectedId} in Byte-Array um
  \item \textbf{HC595::write()} -- Schiebt Bytes zum Schieberegister
\end{enumerate}

\begin{warnbox}[Konsequenz]
LED-Zustände entstehen an genau einer Stelle (\code{setOneHot}). Kein anderer Code
darf LEDs direkt setzen -- das würde die Pipeline umgehen und könnte One-Hot brechen.
\end{warnbox}

% -----------------------------------------------------------------------------
\subsection{One-Hot ist eine systemweite Invariante}
\label{subsec:one-hot}

Bei einem Panel mit 100~LEDs wäre es fatal, wenn versehentlich zwei LEDs gleichzeitig
leuchten oder ein ungültiger Index adressiert wird.

\textbf{Regel:} One-Hot ist ein Vertrag:
\begin{itemize}
  \item \code{selectedId == 0} $\to$ alle LEDs aus
  \item \code{selectedId} $\in [1..N]$ $\to$ genau eine LED an
  \item Andere Zustände sind verboten
\end{itemize}

\paragraph{Implementierung der Invariante}

\begin{lstlisting}[style=arduino,caption={One-Hot Implementierung},label={lst:one-hot}]
void setOneHot(uint8_t selectedId, uint8_t* ledState, size_t ledBytes) {
    // 1. Alle LEDs aus (clear)
    memset(ledState, 0x00, ledBytes);

    // 2. Genau eine LED setzen (wenn selectedId > 0)
    if (selectedId >= 1 && selectedId <= LED_COUNT) {
        uint8_t byteIdx = (selectedId - 1) / 8;
        uint8_t bitIdx  = 7 - ((selectedId - 1) % 8);  // MSB-first
        ledState[byteIdx] |= (1 << bitIdx);
    }

    // 3. Ungueltige Bits im letzten Byte maskieren
    ledState[ledBytes - 1] &= LED_LAST_MASK;
}
\end{lstlisting}

\begin{infobox}[Warum Schritt 3?]
Bei 10~LEDs brauchst du 2~Bytes (16~Bits), aber nur 10 sind gültig. Die Bits 5--0
im zweiten Byte könnten durch Bugs oder Speicherfehler gesetzt sein. Die Maske
garantiert, dass sie niemals an den HC595 gelangen.
\end{infobox}

\paragraph{Beispiel (10 LEDs, MSB-first)}

\begin{table}[H]
\centering
\caption{One-Hot Bit-Mapping für 10 LEDs}
\label{tab:one-hot-mapping}
\begin{tabular}{@{}c l l l@{}}
\toprule
\textbf{selectedId} & \textbf{Byte[0]} & \textbf{Byte[1]} & \textbf{Hex} \\
\midrule
0 (aus) & \code{00000000} & \code{00000000} & \code{0x00 0x00} \\
1       & \code{10000000} & \code{00000000} & \code{0x80 0x00} \\
8       & \code{00000001} & \code{00000000} & \code{0x01 0x00} \\
9       & \code{00000000} & \code{10000000} & \code{0x00 0x80} \\
10      & \code{00000000} & \code{01000000} & \code{0x00 0x40} \\
\bottomrule
\end{tabular}
\end{table}

% -----------------------------------------------------------------------------
\subsection{Ereignisse statt Pegel -- Buttons sind Edge-getrieben}
\label{subsec:edge-driven}

Wenn die Logik auf \glqq Button ist gedrückt\grqq\ (Pegel) reagiert, feuert sie bei
jedem Scan-Zyklus. Bei \SI{10}{\milli\second} Scan-Intervall wären das 100~Events
pro Sekunde, solange du die Taste hältst.

\textbf{Regel:} UI-Aktionen basieren auf Press-Edges (Flanken), nicht auf Pegeln.

\begin{lstlisting}[style=base,numbers=none]
CD4021 (Pegel)      Debounce            Selection
     |                  |                   |
  active-low    ->   pressed edges    ->   reagiert nur
  (0=gedrueckt)     (stabil+Flanke)       auf Edges
\end{lstlisting}

\paragraph{Debounce-Algorithmus (vereinfacht)}

\begin{lstlisting}[style=arduino,caption={Debounce-Implementierung},label={lst:debounce}]
typedef struct {
    uint8_t stable;      // Letzter stabiler Zustand
    uint8_t counter;     // Zaehler fuer Entprellung
} DebounceState;

bool debounce_update(DebounceState* state, uint8_t raw, uint8_t threshold) {
    if (raw == state->stable) {
        state->counter = 0;
        return false;  // Keine Aenderung
    }

    state->counter++;
    if (state->counter >= threshold) {
        bool wasPressed = (state->stable != 0);
        bool isPressed  = (raw != 0);
        state->stable = raw;
        state->counter = 0;

        // Press-Edge: war nicht gedrueckt, ist jetzt gedrueckt
        return (!wasPressed && isPressed);
    }
    return false;
}
\end{lstlisting}

\begin{tipbox}[Konsequenz]
Die Selection-Logik sieht nur \glqq Button 7 wurde gerade gedrückt\grqq, nicht
\glqq Button 7 ist seit \SI{500}{\milli\second} gedrückt\grqq. Das verhindert
Event-Spam und macht das Verhalten vorhersagbar.
\end{tipbox}

% -----------------------------------------------------------------------------
\subsection{Determinismus -- Timing bleibt stabil}
\label{subsec:determinismus}

Im Embedded-Kontext kann ein einzelner \code{Serial.printf()} im Hot-Path das
gesamte Timing zerstören. UART bei \SI{115200}{Baud} braucht etwa \SI{87}{\micro\second}
pro Zeichen -- bei 50~Zeichen sind das \SI{4,3}{\milli\second}, fast die Hälfte
eines \SI{10}{\milli\second}-Scan-Intervalls.

\textbf{Regel:} Der Scan-Pfad läuft ohne Blockierungen. Logging erfolgt nur bei
Zustandsänderungen.

\begin{table}[H]
\centering
\caption{Erlaubte und verbotene Operationen im Hot-Path}
\label{tab:timing-regeln}
\begin{tabularx}{\textwidth}{@{}X X@{}}
\toprule
\textbf{Erlaubt} & \textbf{Verboten} \\
\midrule
Logging bei Selection-Change      & Logging pro Scan-Zyklus \\
\code{vTaskDelayUntil()} (FreeRTOS) & Busy-loops mit \code{delay()} \\
Kurze SPI-Transfers ($\sim$\SI{10}{\micro\second}) & Lange Wartezeiten im Hot-Path \\
\bottomrule
\end{tabularx}
\end{table}

\paragraph{Timing-Budget (10 Kanäle)}

\begin{table}[H]
\centering
\caption{Timing-Budget pro Scan-Zyklus}
\label{tab:timing-budget}
\begin{tabular}{@{}l r@{}}
\toprule
\textbf{Operation} & \textbf{Typische Dauer} \\
\midrule
CD4021 read (2~Bytes)              & $\sim$\SI{20}{\micro\second} \\
Debounce + Selection               & $\sim$\SI{5}{\micro\second} \\
HC595 write (2~Bytes)              & $\sim$\SI{20}{\micro\second} \\
\midrule
\textbf{Gesamt pro Zyklus}         & $\sim$\SI{50}{\micro\second} \\
\textbf{Verfügbar bei \SI{10}{\milli\second} Intervall} & \SI{10000}{\micro\second} \\
\bottomrule
\end{tabular}
\end{table}

\begin{infobox}[Konsequenz]
Bei 10~Kanälen hast du reichlich Reserve. Bei 100~Kanälen (13~Bytes,
$\sim$\SI{130}{\micro\second} Transfer) bleibt das Timing stabil, solange du
keine blockierenden Operationen einbaust.
\end{infobox}

% -----------------------------------------------------------------------------
\subsection{Skalierbarkeit -- Byte-orientiert und maskenbasiert}
\label{subsec:skalierbarkeit}

Code mit hartkodierten Werten wie \code{uint8\_t buttons[2]} oder Schleifen
\code{for (int i = 0; i < 10; i++)} muss bei Skalierung manuell angepasst werden.

\textbf{Regel:} Algorithmen arbeiten mit \code{*\_COUNT} und \code{*\_BYTES},
nicht mit Literalen.

\begin{lstlisting}[style=arduino,caption={Skalierbare vs.\ hartkodierte Implementierung},label={lst:skalierbar}]
// Schlecht: hartkodiert
uint8_t ledState[2];
for (int i = 0; i < 10; i++) { ... }

// Gut: skalierbar
uint8_t ledState[LED_BYTES];
for (int i = 0; i < LED_COUNT; i++) { ... }
\end{lstlisting}

\paragraph{Automatische Berechnung}

\begin{lstlisting}[style=arduino,caption={Automatische Maskenberechnung für 100 Kanäle},label={lst:auto-mask}]
#define BTN_COUNT   100
#define BTN_BYTES   ((BTN_COUNT + 7) / 8)  // = 13

#define LED_COUNT   100
#define LED_BYTES   ((LED_COUNT + 7) / 8)  // = 13

// Maske: Bei 100 LEDs sind im letzten Byte nur Bits 7-4 gueltig (4 LEDs)
// 100 % 8 = 4 -> Bits 7,6,5,4 -> Maske = 0xF0
#define LED_LAST_MASK  (0xFF << (8 - (LED_COUNT % 8))) & 0xFF
\end{lstlisting}

\begin{tipbox}[Konsequenz]
Die Änderung \code{BTN\_COUNT=10} $\to$ \code{BTN\_COUNT=100} erfordert keine
Anpassung in der Logik. Arrays, Schleifen und Masken passen sich automatisch an.
\end{tipbox}

% -----------------------------------------------------------------------------
\subsection{Testbarkeit ist ein Architekturziel}
\label{subsec:testbarkeit}

Wenn Logik und Hardware verwoben sind, kannst du nur am echten Gerät testen.
Jeder Test braucht Flashen, Verkabeln, manuelles Drücken -- langsam und fehleranfällig.

\textbf{Regel:} Logik muss isoliert testbar sein, Hardware-Pfad end-to-end.

\begin{table}[H]
\centering
\caption{Test-Environments in platformio.ini}
\label{tab:test-environments}
\begin{tabularx}{\textwidth}{@{}l l X@{}}
\toprule
\textbf{Environment} & \textbf{Typ} & \textbf{Was wird getestet} \\
\midrule
\code{seeed\_xiao\_esp32s3\_test} & Unit & Nur \code{logic/} (Debounce, Selection, Pipeline) \\
\code{seeed\_xiao\_esp32s3\_e2e}  & E2E  & \code{logic/} + \code{drivers/} + \code{hal/} mit echter Hardware \\
\bottomrule
\end{tabularx}
\end{table}

\paragraph{Unit-Test-Beispiel (Debounce)}

\begin{lstlisting}[style=arduino,caption={Unit-Test für Debounce},label={lst:debounce-test}]
void test_debounce_filters_noise() {
    DebounceState state = {0};

    // Simuliere Prellen: 0 -> 1 -> 0 -> 1 -> 1 -> 1
    TEST_ASSERT_FALSE(debounce_update(&state, 1, 3));  // Zaehler = 1
    TEST_ASSERT_FALSE(debounce_update(&state, 0, 3));  // Reset
    TEST_ASSERT_FALSE(debounce_update(&state, 1, 3));  // Zaehler = 1
    TEST_ASSERT_FALSE(debounce_update(&state, 1, 3));  // Zaehler = 2
    TEST_ASSERT_TRUE(debounce_update(&state, 1, 3));   // Zaehler = 3 -> Edge!
}
\end{lstlisting}

\begin{infobox}[Konsequenz]
Unit-Tests laufen in Sekunden, nicht Minuten. Du kannst Debounce-Schwellenwerte
experimentell optimieren, ohne jedes Mal zu flashen.
\end{infobox}

% -----------------------------------------------------------------------------
\subsection{Änderungen haben Pflichtfolgen}
\label{subsec:pflichtfolgen}

\textbf{Regel:} Bestimmte Änderungen erzwingen bestimmte Checks.

\begin{table}[H]
\centering
\caption{Pflicht-Checks nach Änderungstyp}
\label{tab:pflicht-checks}
\begin{tabularx}{\textwidth}{@{}l X@{}}
\toprule
\textbf{Änderung} & \textbf{Pflicht} \\
\midrule
Pins/SPI/Mapping            & \code{wiring.md} aktualisieren + E2E-Test \\
Debounce/Selection/Pipeline & Unit-Tests (\code{seeed\_xiao\_esp32s3\_test}) \\
Treiber/HAL                 & E2E-Test + Timing-Messung \\
Counts (10 $\to$ 100)       & Unit-Tests + E2E-Test + Timing-Prüfung \\
\bottomrule
\end{tabularx}
\end{table}

\paragraph{Testbefehle}

\begin{lstlisting}[style=shell,caption={Testbefehle für Unit- und E2E-Tests},label={lst:testbefehle}]
# Unit-Tests (Logik)
pio test -e seeed_xiao_esp32s3_test -v

# E2E-Hardware-Test
pio test -e seeed_xiao_esp32s3_e2e -f test_e2e_hw -v
\end{lstlisting}

% -----------------------------------------------------------------------------
\subsection{Minimal-Checkliste}
\label{subsec:checkliste}

Bevor du einen Change committest:

\begin{itemize}
  \item[$\square$] SSOT nur in \code{include/hw/*}
  \item[$\square$] Datenfluss bleibt: read $\to$ update $\to$ write
  \item[$\square$] One-Hot entsteht in genau einem Pfad (clear $\to$ set)
  \item[$\square$] Letztes Byte maskiert (\code{*\_LAST\_MASK})
  \item[$\square$] \code{logic/} ist hardware-unabhängig
  \item[$\square$] Unit-Tests grün (\code{seeed\_xiao\_esp32s3\_test})
  \item[$\square$] Bei HW-relevanten Changes: E2E grün (\code{test\_e2e\_hw})
\end{itemize}

% -----------------------------------------------------------------------------
\subsection{Siehe auch}
\label{subsec:arch-siehe-auch}

Weiterführende Dokumentation:

\begin{itemize}
  \item \code{firmware/docs/architecture.md} -- Schichtenmodell-Details
  \item \code{firmware/docs/overview.md} -- Bit-Mapping (MSB-first)
  \item \code{firmware/docs/wiring.md} -- Hardware-Verdrahtung
  \item \code{apps/bridge/docs/PROTOCOL.md} -- Serial/WebSocket-Protokoll
\end{itemize}

% =============================================================================

\section{Architekturprinzipien -- Robuste Embedded-Software}
\label{sec:architekturprinzipien}

% -----------------------------------------------------------------------------
\subsection{Begriffsrahmen und Zielsetzung}
\label{subsec:arch-begriffsrahmen}

Die Selection-Panel-Firmware steuert 10--100 Buttons und LEDs über Schieberegister
(CD4021, 74HC595). Die Architektur muss drei Anforderungen gleichzeitig erfüllen:
Korrektheit (One-Hot-Invariante), Wartbarkeit (klare Zuständigkeiten), Skalierbarkeit
(10 $\to$ 100 ohne Code-Änderungen in der Logik).

\begin{table}[H]
\centering
\caption{Zentrale Begriffe der Firmware-Architektur}
\label{tab:arch-begriffe}
\begin{tabularx}{\textwidth}{@{}l X@{}}
\toprule
\textbf{Begriff} & \textbf{Definition} \\
\midrule
SSOT       & Single Source of Truth -- eine Konfiguration existiert an genau einer Stelle \\
Invariante & Bedingung, die zu jedem Zeitpunkt gelten muss \\
One-Hot    & Genau ein Bit gesetzt (oder keins), niemals mehrere \\
Pipeline   & Sequenz von Verarbeitungsstufen mit definierten Ein-/Ausgängen \\
Edge       & Zustandsübergang (Flanke), nicht der Zustand selbst \\
HAL        & Hardware Abstraction Layer -- trennt Hardware-Zugriff von Logik \\
\bottomrule
\end{tabularx}
\end{table}

\begin{tipbox}[Ziel]
Stabiler Datenfluss, testbare Logik, keine Seiteneffekte.
\end{tipbox}

% -----------------------------------------------------------------------------
\subsection{SSOT -- Hardware-Konfiguration ist kein \glqq Detail\grqq}
\label{subsec:ssot-prinzip}

Wenn Pins, Zählerstände oder SPI-Parameter an mehreren Stellen definiert sind,
entstehen Inkonsistenzen. Du änderst einen Pin in \code{main.cpp}, vergisst aber
den Treiber -- das System verhält sich unerklärlich.

\textbf{Regel:} Jede Konfiguration existiert an genau einer Stelle. Alle anderen
Dateien referenzieren diese Quelle.

\begin{lstlisting}[style=shell,numbers=none,caption={SSOT-Verzeichnisstruktur},label={lst:ssot-struktur}]
firmware/include/hw/
|-- board_pins.h    # Pin-Definitionen (SSOT)
|-- io_counts.h     # Counts, Masks (SSOT)
`-- spi_params.h    # SPI-Parameter (SSOT)
\end{lstlisting}

\paragraph{Beispiel aus io\_counts.h}

\begin{lstlisting}[style=arduino,caption={SSOT für Counts und Masken},label={lst:io-counts}]
#define BTN_COUNT       10
#define LED_COUNT       10

#define BTN_BYTES       ((BTN_COUNT + 7) / 8)  // = 2
#define LED_BYTES       ((LED_COUNT + 7) / 8)  // = 2

// Maske fuer ungueltige Bits im letzten Byte
#define LED_LAST_MASK   0xC0  // Bei 10 LEDs: nur Bit 7+6 gueltig
\end{lstlisting}

\begin{infobox}[Konsequenz]
Wenn du von 10 auf 100 Kanäle skalierst, änderst du \code{BTN\_COUNT} und
\code{LED\_COUNT} an dieser einen Stelle. Die Makros \code{BTN\_BYTES} und
\code{LED\_LAST\_MASK} berechnen sich automatisch neu.
\end{infobox}

% -----------------------------------------------------------------------------
\subsection{Schichtenmodell -- Abhängigkeiten zeigen nach unten}
\label{subsec:schichtenmodell}

Ohne klare Schichtung entsteht \glqq Spaghetti-Code\grqq: Der Treiber ruft Logik auf,
die Logik greift auf Hardware zu, die App dupliziert Regeln. Änderungen haben
unvorhersehbare Auswirkungen.

\textbf{Regel:} Abhängigkeiten sind gerichtet. Höhere Schichten dürfen niedrigere
nutzen, aber nicht umgekehrt.

\begin{figure}[H]
\centering
\begin{tikzpicture}[
  layer/.style={draw, rectangle, minimum width=11cm, minimum height=0.9cm,
                font=\ttfamily\small, fill=white},
  label/.style={font=\sffamily\footnotesize, text=MutedText}
]
  \node[layer, fill=ArduinoTeal!15] (app) at (0,4) {app/ \hfill Orchestrierung (Loop/Task, Scheduling)};
  \node[layer, fill=ArduinoTeal!10] (logic) at (0,3) {logic/ \hfill Regeln (Debounce, Selection, Pipeline)};
  \node[layer, fill=SuccessGreen!10] (drivers) at (0,2) {drivers/ \hfill IC-spezifisch (CD4021 read, HC595 write)};
  \node[layer, fill=SuccessGreen!15] (hal) at (0,1) {hal/ \hfill SPI-Bus Zugriff (Transfer/Settings)};
  \node[layer, fill=MutedText!10] (hw) at (0,0) {include/hw/* \hfill SSOT (Pins/Counts/SPI-Params)};

  \draw[-{Stealth}, thick] (app.south west) ++(0.3,0) -- ++(0,-0.15);
  \draw[-{Stealth}, thick] (logic.south west) ++(0.3,0) -- ++(0,-0.15);
  \draw[-{Stealth}, thick] (drivers.south west) ++(0.3,0) -- ++(0,-0.15);
  \draw[-{Stealth}, thick] (hal.south west) ++(0.3,0) -- ++(0,-0.15);
\end{tikzpicture}
\caption{Schichtenmodell der Firmware}
\label{fig:schichtenmodell}
\end{figure}

\paragraph{Verzeichnisstruktur}

\begin{lstlisting}[style=shell,numbers=none,caption={Firmware-Verzeichnisstruktur},label={lst:firmware-struktur}]
firmware/src/
|-- hal/
|   `-- spi_bus.cpp         # SPI-Initialisierung
|-- drivers/
|   |-- cd4021.cpp          # Button-Input (liest Bytes)
|   `-- hc595.cpp           # LED-Output (schreibt Bytes)
|-- logic/
|   |-- debounce.cpp        # Entprellung -> Press-Edges
|   |-- selection.cpp       # Selection-Policy -> selectedId
|   `-- io_pipeline.cpp     # Orchestriert Logik-Module
`-- app/
    |-- io_loop.cpp         # Arduino-Variante
    `-- io_task.cpp         # FreeRTOS-Variante
\end{lstlisting}

\begin{table}[H]
\centering
\caption{Schicht-Verantwortlichkeiten}
\label{tab:schicht-regeln}
\begin{tabularx}{\textwidth}{@{}l X X@{}}
\toprule
\textbf{Schicht} & \textbf{Darf} & \textbf{Darf nicht} \\
\midrule
\code{drivers/} & Bytes lesen/schreiben        & Debounce-Entscheidungen treffen \\
\code{logic/}   & IDs und Zustände verarbeiten & Pins kennen, SPI aufrufen \\
\code{app/}     & Schichten orchestrieren      & Regeln duplizieren \\
\bottomrule
\end{tabularx}
\end{table}

\begin{tipbox}[Warum ist das wichtig?]
Die \code{logic/}-Schicht ist hardware-unabhängig. Du kannst sie auf dem Host-PC
unit-testen, ohne ESP32, ohne SPI, ohne echte Hardware.
\end{tipbox}

% -----------------------------------------------------------------------------
\subsection{Datenfluss ist eine Pipeline}
\label{subsec:pipeline}

Wenn IO-Code \glqq verstreut\grqq\ ist -- hier ein LED-Update, dort ein Button-Check --
wird das System unvorhersehbar. Die One-Hot-Invariante kann brechen, weil zwei
Codepfade unabhängig LEDs setzen.

\textbf{Regel:} Der gesamte IO-Zyklus folgt einer klaren Sequenz:

\begin{lstlisting}[style=base,numbers=none]
read(raw) -> normalize/logic(update) -> derive outputs -> write(outputs)
\end{lstlisting}

\begin{figure}[H]
\centering
\begin{tikzpicture}[
  box/.style={draw, rectangle, rounded corners=2pt, minimum height=0.8cm,
              font=\ttfamily\small, fill=white},
  pipeline/.style={draw, rectangle, rounded corners=3pt, minimum width=10cm,
                   minimum height=1.8cm, fill=ArduinoTeal!8},
  arrow/.style={-{Stealth[length=2mm]}, thick}
]
  % CD4021
  \node[box, fill=SuccessGreen!15] (cd4021) at (0,3) {CD4021::read(rawButtons)};

  % Pipeline container
  \node[pipeline] (pipe) at (0,1.2) {};
  \node[font=\sffamily\bfseries\small, anchor=north] at (0,2) {io\_pipeline};

  % Pipeline stages
  \node[box] (debounce) at (-3.5,1.2) {Debounce\\(Edges)};
  \node[box] (guard) at (0,1.2) {Press-Guard\\(optional)};
  \node[box] (selection) at (3.5,1.2) {Selection\\(Policy)};

  \draw[arrow] (debounce) -- (guard);
  \draw[arrow] (guard) -- (selection);

  % selectedId
  \node[font=\ttfamily\small] (selid) at (0,-0.3) {selectedId};

  % setOneHot
  \node[box, fill=ArduinoTeal!15] (onehot) at (0,-1.3) {setOneHot(selectedId) $\to$ ledState[]};

  % HC595
  \node[box, fill=SuccessGreen!15] (hc595) at (0,-2.5) {HC595::write(ledState)};

  % Arrows
  \draw[arrow] (cd4021) -- (pipe.north);
  \draw[arrow] (pipe.south) -- (selid);
  \draw[arrow] (selid) -- (onehot);
  \draw[arrow] (onehot) -- (hc595);
\end{tikzpicture}
\caption{Datenfluss-Pipeline im Selection Panel}
\label{fig:pipeline}
\end{figure}

\paragraph{Was passiert in jeder Stufe?}

\begin{enumerate}
  \item \textbf{CD4021::read()} -- Liest rohe Bytes vom Schieberegister (active-low, 0~=~gedrückt)
  \item \textbf{Debounce} -- Filtert Prellen, erzeugt stabile Press-Edges
  \item \textbf{Press-Guard} -- Verhindert versehentliche Doppel-Presses (optional)
  \item \textbf{Selection} -- Wendet Policy an: \glqq Welche LED soll leuchten?\grqq
  \item \textbf{setOneHot()} -- Wandelt \code{selectedId} in Byte-Array um
  \item \textbf{HC595::write()} -- Schiebt Bytes zum Schieberegister
\end{enumerate}

\begin{warnbox}[Konsequenz]
LED-Zustände entstehen an genau einer Stelle (\code{setOneHot}). Kein anderer Code
darf LEDs direkt setzen -- das würde die Pipeline umgehen und könnte One-Hot brechen.
\end{warnbox}

% -----------------------------------------------------------------------------
\subsection{One-Hot ist eine systemweite Invariante}
\label{subsec:one-hot}

Bei einem Panel mit 100~LEDs wäre es fatal, wenn versehentlich zwei LEDs gleichzeitig
leuchten oder ein ungültiger Index adressiert wird.

\textbf{Regel:} One-Hot ist ein Vertrag:
\begin{itemize}
  \item \code{selectedId == 0} $\to$ alle LEDs aus
  \item \code{selectedId} $\in [1..N]$ $\to$ genau eine LED an
  \item Andere Zustände sind verboten
\end{itemize}

\paragraph{Implementierung der Invariante}

\begin{lstlisting}[style=arduino,caption={One-Hot Implementierung},label={lst:one-hot}]
void setOneHot(uint8_t selectedId, uint8_t* ledState, size_t ledBytes) {
    // 1. Alle LEDs aus (clear)
    memset(ledState, 0x00, ledBytes);

    // 2. Genau eine LED setzen (wenn selectedId > 0)
    if (selectedId >= 1 && selectedId <= LED_COUNT) {
        uint8_t byteIdx = (selectedId - 1) / 8;
        uint8_t bitIdx  = 7 - ((selectedId - 1) % 8);  // MSB-first
        ledState[byteIdx] |= (1 << bitIdx);
    }

    // 3. Ungueltige Bits im letzten Byte maskieren
    ledState[ledBytes - 1] &= LED_LAST_MASK;
}
\end{lstlisting}

\begin{infobox}[Warum Schritt 3?]
Bei 10~LEDs brauchst du 2~Bytes (16~Bits), aber nur 10 sind gültig. Die Bits 5--0
im zweiten Byte könnten durch Bugs oder Speicherfehler gesetzt sein. Die Maske
garantiert, dass sie niemals an den HC595 gelangen.
\end{infobox}

\paragraph{Beispiel (10 LEDs, MSB-first)}

\begin{table}[H]
\centering
\caption{One-Hot Bit-Mapping für 10 LEDs}
\label{tab:one-hot-mapping}
\begin{tabular}{@{}c l l l@{}}
\toprule
\textbf{selectedId} & \textbf{Byte[0]} & \textbf{Byte[1]} & \textbf{Hex} \\
\midrule
0 (aus) & \code{00000000} & \code{00000000} & \code{0x00 0x00} \\
1       & \code{10000000} & \code{00000000} & \code{0x80 0x00} \\
8       & \code{00000001} & \code{00000000} & \code{0x01 0x00} \\
9       & \code{00000000} & \code{10000000} & \code{0x00 0x80} \\
10      & \code{00000000} & \code{01000000} & \code{0x00 0x40} \\
\bottomrule
\end{tabular}
\end{table}

% -----------------------------------------------------------------------------
\subsection{Ereignisse statt Pegel -- Buttons sind Edge-getrieben}
\label{subsec:edge-driven}

Wenn die Logik auf \glqq Button ist gedrückt\grqq\ (Pegel) reagiert, feuert sie bei
jedem Scan-Zyklus. Bei \SI{10}{\milli\second} Scan-Intervall wären das 100~Events
pro Sekunde, solange du die Taste hältst.

\textbf{Regel:} UI-Aktionen basieren auf Press-Edges (Flanken), nicht auf Pegeln.

\begin{lstlisting}[style=base,numbers=none]
CD4021 (Pegel)      Debounce            Selection
     |                  |                   |
  active-low    ->   pressed edges    ->   reagiert nur
  (0=gedrueckt)     (stabil+Flanke)       auf Edges
\end{lstlisting}

\paragraph{Debounce-Algorithmus (vereinfacht)}

\begin{lstlisting}[style=arduino,caption={Debounce-Implementierung},label={lst:debounce}]
typedef struct {
    uint8_t stable;      // Letzter stabiler Zustand
    uint8_t counter;     // Zaehler fuer Entprellung
} DebounceState;

bool debounce_update(DebounceState* state, uint8_t raw, uint8_t threshold) {
    if (raw == state->stable) {
        state->counter = 0;
        return false;  // Keine Aenderung
    }

    state->counter++;
    if (state->counter >= threshold) {
        bool wasPressed = (state->stable != 0);
        bool isPressed  = (raw != 0);
        state->stable = raw;
        state->counter = 0;

        // Press-Edge: war nicht gedrueckt, ist jetzt gedrueckt
        return (!wasPressed && isPressed);
    }
    return false;
}
\end{lstlisting}

\begin{tipbox}[Konsequenz]
Die Selection-Logik sieht nur \glqq Button 7 wurde gerade gedrückt\grqq, nicht
\glqq Button 7 ist seit \SI{500}{\milli\second} gedrückt\grqq. Das verhindert
Event-Spam und macht das Verhalten vorhersagbar.
\end{tipbox}

% -----------------------------------------------------------------------------
\subsection{Determinismus -- Timing bleibt stabil}
\label{subsec:determinismus}

Im Embedded-Kontext kann ein einzelner \code{Serial.printf()} im Hot-Path das
gesamte Timing zerstören. UART bei \SI{115200}{Baud} braucht etwa \SI{87}{\micro\second}
pro Zeichen -- bei 50~Zeichen sind das \SI{4,3}{\milli\second}, fast die Hälfte
eines \SI{10}{\milli\second}-Scan-Intervalls.

\textbf{Regel:} Der Scan-Pfad läuft ohne Blockierungen. Logging erfolgt nur bei
Zustandsänderungen.

\begin{table}[H]
\centering
\caption{Erlaubte und verbotene Operationen im Hot-Path}
\label{tab:timing-regeln}
\begin{tabularx}{\textwidth}{@{}X X@{}}
\toprule
\textbf{Erlaubt} & \textbf{Verboten} \\
\midrule
Logging bei Selection-Change      & Logging pro Scan-Zyklus \\
\code{vTaskDelayUntil()} (FreeRTOS) & Busy-loops mit \code{delay()} \\
Kurze SPI-Transfers ($\sim$\SI{10}{\micro\second}) & Lange Wartezeiten im Hot-Path \\
\bottomrule
\end{tabularx}
\end{table}

\paragraph{Timing-Budget (10 Kanäle)}

\begin{table}[H]
\centering
\caption{Timing-Budget pro Scan-Zyklus}
\label{tab:timing-budget}
\begin{tabular}{@{}l r@{}}
\toprule
\textbf{Operation} & \textbf{Typische Dauer} \\
\midrule
CD4021 read (2~Bytes)              & $\sim$\SI{20}{\micro\second} \\
Debounce + Selection               & $\sim$\SI{5}{\micro\second} \\
HC595 write (2~Bytes)              & $\sim$\SI{20}{\micro\second} \\
\midrule
\textbf{Gesamt pro Zyklus}         & $\sim$\SI{50}{\micro\second} \\
\textbf{Verfügbar bei \SI{10}{\milli\second} Intervall} & \SI{10000}{\micro\second} \\
\bottomrule
\end{tabular}
\end{table}

\begin{infobox}[Konsequenz]
Bei 10~Kanälen hast du reichlich Reserve. Bei 100~Kanälen (13~Bytes,
$\sim$\SI{130}{\micro\second} Transfer) bleibt das Timing stabil, solange du
keine blockierenden Operationen einbaust.
\end{infobox}

% -----------------------------------------------------------------------------
\subsection{Skalierbarkeit -- Byte-orientiert und maskenbasiert}
\label{subsec:skalierbarkeit}

Code mit hartkodierten Werten wie \code{uint8\_t buttons[2]} oder Schleifen
\code{for (int i = 0; i < 10; i++)} muss bei Skalierung manuell angepasst werden.

\textbf{Regel:} Algorithmen arbeiten mit \code{*\_COUNT} und \code{*\_BYTES},
nicht mit Literalen.

\begin{lstlisting}[style=arduino,caption={Skalierbare vs.\ hartkodierte Implementierung},label={lst:skalierbar}]
// Schlecht: hartkodiert
uint8_t ledState[2];
for (int i = 0; i < 10; i++) { ... }

// Gut: skalierbar
uint8_t ledState[LED_BYTES];
for (int i = 0; i < LED_COUNT; i++) { ... }
\end{lstlisting}

\paragraph{Automatische Berechnung}

\begin{lstlisting}[style=arduino,caption={Automatische Maskenberechnung für 100 Kanäle},label={lst:auto-mask}]
#define BTN_COUNT   100
#define BTN_BYTES   ((BTN_COUNT + 7) / 8)  // = 13

#define LED_COUNT   100
#define LED_BYTES   ((LED_COUNT + 7) / 8)  // = 13

// Maske: Bei 100 LEDs sind im letzten Byte nur Bits 7-4 gueltig (4 LEDs)
// 100 % 8 = 4 -> Bits 7,6,5,4 -> Maske = 0xF0
#define LED_LAST_MASK  (0xFF << (8 - (LED_COUNT % 8))) & 0xFF
\end{lstlisting}

\begin{tipbox}[Konsequenz]
Die Änderung \code{BTN\_COUNT=10} $\to$ \code{BTN\_COUNT=100} erfordert keine
Anpassung in der Logik. Arrays, Schleifen und Masken passen sich automatisch an.
\end{tipbox}

% -----------------------------------------------------------------------------
\subsection{Testbarkeit ist ein Architekturziel}
\label{subsec:testbarkeit}

Wenn Logik und Hardware verwoben sind, kannst du nur am echten Gerät testen.
Jeder Test braucht Flashen, Verkabeln, manuelles Drücken -- langsam und fehleranfällig.

\textbf{Regel:} Logik muss isoliert testbar sein, Hardware-Pfad end-to-end.

\begin{table}[H]
\centering
\caption{Test-Environments in platformio.ini}
\label{tab:test-environments}
\begin{tabularx}{\textwidth}{@{}l l X@{}}
\toprule
\textbf{Environment} & \textbf{Typ} & \textbf{Was wird getestet} \\
\midrule
\code{seeed\_xiao\_esp32s3\_test} & Unit & Nur \code{logic/} (Debounce, Selection, Pipeline) \\
\code{seeed\_xiao\_esp32s3\_e2e}  & E2E  & \code{logic/} + \code{drivers/} + \code{hal/} mit echter Hardware \\
\bottomrule
\end{tabularx}
\end{table}

\paragraph{Unit-Test-Beispiel (Debounce)}

\begin{lstlisting}[style=arduino,caption={Unit-Test für Debounce},label={lst:debounce-test}]
void test_debounce_filters_noise() {
    DebounceState state = {0};

    // Simuliere Prellen: 0 -> 1 -> 0 -> 1 -> 1 -> 1
    TEST_ASSERT_FALSE(debounce_update(&state, 1, 3));  // Zaehler = 1
    TEST_ASSERT_FALSE(debounce_update(&state, 0, 3));  // Reset
    TEST_ASSERT_FALSE(debounce_update(&state, 1, 3));  // Zaehler = 1
    TEST_ASSERT_FALSE(debounce_update(&state, 1, 3));  // Zaehler = 2
    TEST_ASSERT_TRUE(debounce_update(&state, 1, 3));   // Zaehler = 3 -> Edge!
}
\end{lstlisting}

\begin{infobox}[Konsequenz]
Unit-Tests laufen in Sekunden, nicht Minuten. Du kannst Debounce-Schwellenwerte
experimentell optimieren, ohne jedes Mal zu flashen.
\end{infobox}

% -----------------------------------------------------------------------------
\subsection{Änderungen haben Pflichtfolgen}
\label{subsec:pflichtfolgen}

\textbf{Regel:} Bestimmte Änderungen erzwingen bestimmte Checks.

\begin{table}[H]
\centering
\caption{Pflicht-Checks nach Änderungstyp}
\label{tab:pflicht-checks}
\begin{tabularx}{\textwidth}{@{}l X@{}}
\toprule
\textbf{Änderung} & \textbf{Pflicht} \\
\midrule
Pins/SPI/Mapping            & \code{wiring.md} aktualisieren + E2E-Test \\
Debounce/Selection/Pipeline & Unit-Tests (\code{seeed\_xiao\_esp32s3\_test}) \\
Treiber/HAL                 & E2E-Test + Timing-Messung \\
Counts (10 $\to$ 100)       & Unit-Tests + E2E-Test + Timing-Prüfung \\
\bottomrule
\end{tabularx}
\end{table}

\paragraph{Testbefehle}

\begin{lstlisting}[style=shell,caption={Testbefehle für Unit- und E2E-Tests},label={lst:testbefehle}]
# Unit-Tests (Logik)
pio test -e seeed_xiao_esp32s3_test -v

# E2E-Hardware-Test
pio test -e seeed_xiao_esp32s3_e2e -f test_e2e_hw -v
\end{lstlisting}

% -----------------------------------------------------------------------------
\subsection{Minimal-Checkliste}
\label{subsec:checkliste}

Bevor du einen Change committest:

\begin{itemize}
  \item[$\square$] SSOT nur in \code{include/hw/*}
  \item[$\square$] Datenfluss bleibt: read $\to$ update $\to$ write
  \item[$\square$] One-Hot entsteht in genau einem Pfad (clear $\to$ set)
  \item[$\square$] Letztes Byte maskiert (\code{*\_LAST\_MASK})
  \item[$\square$] \code{logic/} ist hardware-unabhängig
  \item[$\square$] Unit-Tests grün (\code{seeed\_xiao\_esp32s3\_test})
  \item[$\square$] Bei HW-relevanten Changes: E2E grün (\code{test\_e2e\_hw})
\end{itemize}

% -----------------------------------------------------------------------------
\subsection{Siehe auch}
\label{subsec:arch-siehe-auch}

Weiterführende Dokumentation:

\begin{itemize}
  \item \code{firmware/docs/architecture.md} -- Schichtenmodell-Details
  \item \code{firmware/docs/overview.md} -- Bit-Mapping (MSB-first)
  \item \code{firmware/docs/wiring.md} -- Hardware-Verdrahtung
  \item \code{apps/bridge/docs/PROTOCOL.md} -- Serial/WebSocket-Protokoll
\end{itemize}

% =============================================================================

\section{Architekturprinzipien -- Robuste Embedded-Software}
\label{sec:architekturprinzipien}

% -----------------------------------------------------------------------------
\subsection{Begriffsrahmen und Zielsetzung}
\label{subsec:arch-begriffsrahmen}

Die Selection-Panel-Firmware steuert 10--100 Buttons und LEDs über Schieberegister
(CD4021, 74HC595). Die Architektur muss drei Anforderungen gleichzeitig erfüllen:
Korrektheit (One-Hot-Invariante), Wartbarkeit (klare Zuständigkeiten), Skalierbarkeit
(10 $\to$ 100 ohne Code-Änderungen in der Logik).

\begin{table}[H]
\centering
\caption{Zentrale Begriffe der Firmware-Architektur}
\label{tab:arch-begriffe}
\begin{tabularx}{\textwidth}{@{}l X@{}}
\toprule
\textbf{Begriff} & \textbf{Definition} \\
\midrule
SSOT       & Single Source of Truth -- eine Konfiguration existiert an genau einer Stelle \\
Invariante & Bedingung, die zu jedem Zeitpunkt gelten muss \\
One-Hot    & Genau ein Bit gesetzt (oder keins), niemals mehrere \\
Pipeline   & Sequenz von Verarbeitungsstufen mit definierten Ein-/Ausgängen \\
Edge       & Zustandsübergang (Flanke), nicht der Zustand selbst \\
HAL        & Hardware Abstraction Layer -- trennt Hardware-Zugriff von Logik \\
\bottomrule
\end{tabularx}
\end{table}

\begin{tipbox}[Ziel]
Stabiler Datenfluss, testbare Logik, keine Seiteneffekte.
\end{tipbox}

% -----------------------------------------------------------------------------
\subsection{SSOT -- Hardware-Konfiguration ist kein \glqq Detail\grqq}
\label{subsec:ssot-prinzip}

Wenn Pins, Zählerstände oder SPI-Parameter an mehreren Stellen definiert sind,
entstehen Inkonsistenzen. Du änderst einen Pin in \code{main.cpp}, vergisst aber
den Treiber -- das System verhält sich unerklärlich.

\textbf{Regel:} Jede Konfiguration existiert an genau einer Stelle. Alle anderen
Dateien referenzieren diese Quelle.

\begin{lstlisting}[style=shell,numbers=none,caption={SSOT-Verzeichnisstruktur},label={lst:ssot-struktur}]
firmware/include/hw/
|-- board_pins.h    # Pin-Definitionen (SSOT)
|-- io_counts.h     # Counts, Masks (SSOT)
`-- spi_params.h    # SPI-Parameter (SSOT)
\end{lstlisting}

\paragraph{Beispiel aus io\_counts.h}

\begin{lstlisting}[style=arduino,caption={SSOT für Counts und Masken},label={lst:io-counts}]
#define BTN_COUNT       10
#define LED_COUNT       10

#define BTN_BYTES       ((BTN_COUNT + 7) / 8)  // = 2
#define LED_BYTES       ((LED_COUNT + 7) / 8)  // = 2

// Maske fuer ungueltige Bits im letzten Byte
#define LED_LAST_MASK   0xC0  // Bei 10 LEDs: nur Bit 7+6 gueltig
\end{lstlisting}

\begin{infobox}[Konsequenz]
Wenn du von 10 auf 100 Kanäle skalierst, änderst du \code{BTN\_COUNT} und
\code{LED\_COUNT} an dieser einen Stelle. Die Makros \code{BTN\_BYTES} und
\code{LED\_LAST\_MASK} berechnen sich automatisch neu.
\end{infobox}

% -----------------------------------------------------------------------------
\subsection{Schichtenmodell -- Abhängigkeiten zeigen nach unten}
\label{subsec:schichtenmodell}

Ohne klare Schichtung entsteht \glqq Spaghetti-Code\grqq: Der Treiber ruft Logik auf,
die Logik greift auf Hardware zu, die App dupliziert Regeln. Änderungen haben
unvorhersehbare Auswirkungen.

\textbf{Regel:} Abhängigkeiten sind gerichtet. Höhere Schichten dürfen niedrigere
nutzen, aber nicht umgekehrt.

\begin{figure}[H]
\centering
\begin{tikzpicture}[
  layer/.style={draw, rectangle, minimum width=11cm, minimum height=0.9cm,
                font=\ttfamily\small, fill=white},
  label/.style={font=\sffamily\footnotesize, text=MutedText}
]
  \node[layer, fill=ArduinoTeal!15] (app) at (0,4) {app/ \hfill Orchestrierung (Loop/Task, Scheduling)};
  \node[layer, fill=ArduinoTeal!10] (logic) at (0,3) {logic/ \hfill Regeln (Debounce, Selection, Pipeline)};
  \node[layer, fill=SuccessGreen!10] (drivers) at (0,2) {drivers/ \hfill IC-spezifisch (CD4021 read, HC595 write)};
  \node[layer, fill=SuccessGreen!15] (hal) at (0,1) {hal/ \hfill SPI-Bus Zugriff (Transfer/Settings)};
  \node[layer, fill=MutedText!10] (hw) at (0,0) {include/hw/* \hfill SSOT (Pins/Counts/SPI-Params)};

  \draw[-{Stealth}, thick] (app.south west) ++(0.3,0) -- ++(0,-0.15);
  \draw[-{Stealth}, thick] (logic.south west) ++(0.3,0) -- ++(0,-0.15);
  \draw[-{Stealth}, thick] (drivers.south west) ++(0.3,0) -- ++(0,-0.15);
  \draw[-{Stealth}, thick] (hal.south west) ++(0.3,0) -- ++(0,-0.15);
\end{tikzpicture}
\caption{Schichtenmodell der Firmware}
\label{fig:schichtenmodell}
\end{figure}

\paragraph{Verzeichnisstruktur}

\begin{lstlisting}[style=shell,numbers=none,caption={Firmware-Verzeichnisstruktur},label={lst:firmware-struktur}]
firmware/src/
|-- hal/
|   `-- spi_bus.cpp         # SPI-Initialisierung
|-- drivers/
|   |-- cd4021.cpp          # Button-Input (liest Bytes)
|   `-- hc595.cpp           # LED-Output (schreibt Bytes)
|-- logic/
|   |-- debounce.cpp        # Entprellung -> Press-Edges
|   |-- selection.cpp       # Selection-Policy -> selectedId
|   `-- io_pipeline.cpp     # Orchestriert Logik-Module
`-- app/
    |-- io_loop.cpp         # Arduino-Variante
    `-- io_task.cpp         # FreeRTOS-Variante
\end{lstlisting}

\begin{table}[H]
\centering
\caption{Schicht-Verantwortlichkeiten}
\label{tab:schicht-regeln}
\begin{tabularx}{\textwidth}{@{}l X X@{}}
\toprule
\textbf{Schicht} & \textbf{Darf} & \textbf{Darf nicht} \\
\midrule
\code{drivers/} & Bytes lesen/schreiben        & Debounce-Entscheidungen treffen \\
\code{logic/}   & IDs und Zustände verarbeiten & Pins kennen, SPI aufrufen \\
\code{app/}     & Schichten orchestrieren      & Regeln duplizieren \\
\bottomrule
\end{tabularx}
\end{table}

\begin{tipbox}[Warum ist das wichtig?]
Die \code{logic/}-Schicht ist hardware-unabhängig. Du kannst sie auf dem Host-PC
unit-testen, ohne ESP32, ohne SPI, ohne echte Hardware.
\end{tipbox}

% -----------------------------------------------------------------------------
\subsection{Datenfluss ist eine Pipeline}
\label{subsec:pipeline}

Wenn IO-Code \glqq verstreut\grqq\ ist -- hier ein LED-Update, dort ein Button-Check --
wird das System unvorhersehbar. Die One-Hot-Invariante kann brechen, weil zwei
Codepfade unabhängig LEDs setzen.

\textbf{Regel:} Der gesamte IO-Zyklus folgt einer klaren Sequenz:

\begin{lstlisting}[style=base,numbers=none]
read(raw) -> normalize/logic(update) -> derive outputs -> write(outputs)
\end{lstlisting}

\begin{figure}[H]
\centering
\begin{tikzpicture}[
  box/.style={draw, rectangle, rounded corners=2pt, minimum height=0.8cm,
              font=\ttfamily\small, fill=white},
  pipeline/.style={draw, rectangle, rounded corners=3pt, minimum width=10cm,
                   minimum height=1.8cm, fill=ArduinoTeal!8},
  arrow/.style={-{Stealth[length=2mm]}, thick}
]
  % CD4021
  \node[box, fill=SuccessGreen!15] (cd4021) at (0,3) {CD4021::read(rawButtons)};

  % Pipeline container
  \node[pipeline] (pipe) at (0,1.2) {};
  \node[font=\sffamily\bfseries\small, anchor=north] at (0,2) {io\_pipeline};

  % Pipeline stages
  \node[box] (debounce) at (-3.5,1.2) {Debounce\\(Edges)};
  \node[box] (guard) at (0,1.2) {Press-Guard\\(optional)};
  \node[box] (selection) at (3.5,1.2) {Selection\\(Policy)};

  \draw[arrow] (debounce) -- (guard);
  \draw[arrow] (guard) -- (selection);

  % selectedId
  \node[font=\ttfamily\small] (selid) at (0,-0.3) {selectedId};

  % setOneHot
  \node[box, fill=ArduinoTeal!15] (onehot) at (0,-1.3) {setOneHot(selectedId) $\to$ ledState[]};

  % HC595
  \node[box, fill=SuccessGreen!15] (hc595) at (0,-2.5) {HC595::write(ledState)};

  % Arrows
  \draw[arrow] (cd4021) -- (pipe.north);
  \draw[arrow] (pipe.south) -- (selid);
  \draw[arrow] (selid) -- (onehot);
  \draw[arrow] (onehot) -- (hc595);
\end{tikzpicture}
\caption{Datenfluss-Pipeline im Selection Panel}
\label{fig:pipeline}
\end{figure}

\paragraph{Was passiert in jeder Stufe?}

\begin{enumerate}
  \item \textbf{CD4021::read()} -- Liest rohe Bytes vom Schieberegister (active-low, 0~=~gedrückt)
  \item \textbf{Debounce} -- Filtert Prellen, erzeugt stabile Press-Edges
  \item \textbf{Press-Guard} -- Verhindert versehentliche Doppel-Presses (optional)
  \item \textbf{Selection} -- Wendet Policy an: \glqq Welche LED soll leuchten?\grqq
  \item \textbf{setOneHot()} -- Wandelt \code{selectedId} in Byte-Array um
  \item \textbf{HC595::write()} -- Schiebt Bytes zum Schieberegister
\end{enumerate}

\begin{warnbox}[Konsequenz]
LED-Zustände entstehen an genau einer Stelle (\code{setOneHot}). Kein anderer Code
darf LEDs direkt setzen -- das würde die Pipeline umgehen und könnte One-Hot brechen.
\end{warnbox}

% -----------------------------------------------------------------------------
\subsection{One-Hot ist eine systemweite Invariante}
\label{subsec:one-hot}

Bei einem Panel mit 100~LEDs wäre es fatal, wenn versehentlich zwei LEDs gleichzeitig
leuchten oder ein ungültiger Index adressiert wird.

\textbf{Regel:} One-Hot ist ein Vertrag:
\begin{itemize}
  \item \code{selectedId == 0} $\to$ alle LEDs aus
  \item \code{selectedId} $\in [1..N]$ $\to$ genau eine LED an
  \item Andere Zustände sind verboten
\end{itemize}

\paragraph{Implementierung der Invariante}

\begin{lstlisting}[style=arduino,caption={One-Hot Implementierung},label={lst:one-hot}]
void setOneHot(uint8_t selectedId, uint8_t* ledState, size_t ledBytes) {
    // 1. Alle LEDs aus (clear)
    memset(ledState, 0x00, ledBytes);

    // 2. Genau eine LED setzen (wenn selectedId > 0)
    if (selectedId >= 1 && selectedId <= LED_COUNT) {
        uint8_t byteIdx = (selectedId - 1) / 8;
        uint8_t bitIdx  = 7 - ((selectedId - 1) % 8);  // MSB-first
        ledState[byteIdx] |= (1 << bitIdx);
    }

    // 3. Ungueltige Bits im letzten Byte maskieren
    ledState[ledBytes - 1] &= LED_LAST_MASK;
}
\end{lstlisting}

\begin{infobox}[Warum Schritt 3?]
Bei 10~LEDs brauchst du 2~Bytes (16~Bits), aber nur 10 sind gültig. Die Bits 5--0
im zweiten Byte könnten durch Bugs oder Speicherfehler gesetzt sein. Die Maske
garantiert, dass sie niemals an den HC595 gelangen.
\end{infobox}

\paragraph{Beispiel (10 LEDs, MSB-first)}

\begin{table}[H]
\centering
\caption{One-Hot Bit-Mapping für 10 LEDs}
\label{tab:one-hot-mapping}
\begin{tabular}{@{}c l l l@{}}
\toprule
\textbf{selectedId} & \textbf{Byte[0]} & \textbf{Byte[1]} & \textbf{Hex} \\
\midrule
0 (aus) & \code{00000000} & \code{00000000} & \code{0x00 0x00} \\
1       & \code{10000000} & \code{00000000} & \code{0x80 0x00} \\
8       & \code{00000001} & \code{00000000} & \code{0x01 0x00} \\
9       & \code{00000000} & \code{10000000} & \code{0x00 0x80} \\
10      & \code{00000000} & \code{01000000} & \code{0x00 0x40} \\
\bottomrule
\end{tabular}
\end{table}

% -----------------------------------------------------------------------------
\subsection{Ereignisse statt Pegel -- Buttons sind Edge-getrieben}
\label{subsec:edge-driven}

Wenn die Logik auf \glqq Button ist gedrückt\grqq\ (Pegel) reagiert, feuert sie bei
jedem Scan-Zyklus. Bei \SI{10}{\milli\second} Scan-Intervall wären das 100~Events
pro Sekunde, solange du die Taste hältst.

\textbf{Regel:} UI-Aktionen basieren auf Press-Edges (Flanken), nicht auf Pegeln.

\begin{lstlisting}[style=base,numbers=none]
CD4021 (Pegel)      Debounce            Selection
     |                  |                   |
  active-low    ->   pressed edges    ->   reagiert nur
  (0=gedrueckt)     (stabil+Flanke)       auf Edges
\end{lstlisting}

\paragraph{Debounce-Algorithmus (vereinfacht)}

\begin{lstlisting}[style=arduino,caption={Debounce-Implementierung},label={lst:debounce}]
typedef struct {
    uint8_t stable;      // Letzter stabiler Zustand
    uint8_t counter;     // Zaehler fuer Entprellung
} DebounceState;

bool debounce_update(DebounceState* state, uint8_t raw, uint8_t threshold) {
    if (raw == state->stable) {
        state->counter = 0;
        return false;  // Keine Aenderung
    }

    state->counter++;
    if (state->counter >= threshold) {
        bool wasPressed = (state->stable != 0);
        bool isPressed  = (raw != 0);
        state->stable = raw;
        state->counter = 0;

        // Press-Edge: war nicht gedrueckt, ist jetzt gedrueckt
        return (!wasPressed && isPressed);
    }
    return false;
}
\end{lstlisting}

\begin{tipbox}[Konsequenz]
Die Selection-Logik sieht nur \glqq Button 7 wurde gerade gedrückt\grqq, nicht
\glqq Button 7 ist seit \SI{500}{\milli\second} gedrückt\grqq. Das verhindert
Event-Spam und macht das Verhalten vorhersagbar.
\end{tipbox}

% -----------------------------------------------------------------------------
\subsection{Determinismus -- Timing bleibt stabil}
\label{subsec:determinismus}

Im Embedded-Kontext kann ein einzelner \code{Serial.printf()} im Hot-Path das
gesamte Timing zerstören. UART bei \SI{115200}{Baud} braucht etwa \SI{87}{\micro\second}
pro Zeichen -- bei 50~Zeichen sind das \SI{4,3}{\milli\second}, fast die Hälfte
eines \SI{10}{\milli\second}-Scan-Intervalls.

\textbf{Regel:} Der Scan-Pfad läuft ohne Blockierungen. Logging erfolgt nur bei
Zustandsänderungen.

\begin{table}[H]
\centering
\caption{Erlaubte und verbotene Operationen im Hot-Path}
\label{tab:timing-regeln}
\begin{tabularx}{\textwidth}{@{}X X@{}}
\toprule
\textbf{Erlaubt} & \textbf{Verboten} \\
\midrule
Logging bei Selection-Change      & Logging pro Scan-Zyklus \\
\code{vTaskDelayUntil()} (FreeRTOS) & Busy-loops mit \code{delay()} \\
Kurze SPI-Transfers ($\sim$\SI{10}{\micro\second}) & Lange Wartezeiten im Hot-Path \\
\bottomrule
\end{tabularx}
\end{table}

\paragraph{Timing-Budget (10 Kanäle)}

\begin{table}[H]
\centering
\caption{Timing-Budget pro Scan-Zyklus}
\label{tab:timing-budget}
\begin{tabular}{@{}l r@{}}
\toprule
\textbf{Operation} & \textbf{Typische Dauer} \\
\midrule
CD4021 read (2~Bytes)              & $\sim$\SI{20}{\micro\second} \\
Debounce + Selection               & $\sim$\SI{5}{\micro\second} \\
HC595 write (2~Bytes)              & $\sim$\SI{20}{\micro\second} \\
\midrule
\textbf{Gesamt pro Zyklus}         & $\sim$\SI{50}{\micro\second} \\
\textbf{Verfügbar bei \SI{10}{\milli\second} Intervall} & \SI{10000}{\micro\second} \\
\bottomrule
\end{tabular}
\end{table}

\begin{infobox}[Konsequenz]
Bei 10~Kanälen hast du reichlich Reserve. Bei 100~Kanälen (13~Bytes,
$\sim$\SI{130}{\micro\second} Transfer) bleibt das Timing stabil, solange du
keine blockierenden Operationen einbaust.
\end{infobox}

% -----------------------------------------------------------------------------
\subsection{Skalierbarkeit -- Byte-orientiert und maskenbasiert}
\label{subsec:skalierbarkeit}

Code mit hartkodierten Werten wie \code{uint8\_t buttons[2]} oder Schleifen
\code{for (int i = 0; i < 10; i++)} muss bei Skalierung manuell angepasst werden.

\textbf{Regel:} Algorithmen arbeiten mit \code{*\_COUNT} und \code{*\_BYTES},
nicht mit Literalen.

\begin{lstlisting}[style=arduino,caption={Skalierbare vs.\ hartkodierte Implementierung},label={lst:skalierbar}]
// Schlecht: hartkodiert
uint8_t ledState[2];
for (int i = 0; i < 10; i++) { ... }

// Gut: skalierbar
uint8_t ledState[LED_BYTES];
for (int i = 0; i < LED_COUNT; i++) { ... }
\end{lstlisting}

\paragraph{Automatische Berechnung}

\begin{lstlisting}[style=arduino,caption={Automatische Maskenberechnung für 100 Kanäle},label={lst:auto-mask}]
#define BTN_COUNT   100
#define BTN_BYTES   ((BTN_COUNT + 7) / 8)  // = 13

#define LED_COUNT   100
#define LED_BYTES   ((LED_COUNT + 7) / 8)  // = 13

// Maske: Bei 100 LEDs sind im letzten Byte nur Bits 7-4 gueltig (4 LEDs)
// 100 % 8 = 4 -> Bits 7,6,5,4 -> Maske = 0xF0
#define LED_LAST_MASK  (0xFF << (8 - (LED_COUNT % 8))) & 0xFF
\end{lstlisting}

\begin{tipbox}[Konsequenz]
Die Änderung \code{BTN\_COUNT=10} $\to$ \code{BTN\_COUNT=100} erfordert keine
Anpassung in der Logik. Arrays, Schleifen und Masken passen sich automatisch an.
\end{tipbox}

% -----------------------------------------------------------------------------
\subsection{Testbarkeit ist ein Architekturziel}
\label{subsec:testbarkeit}

Wenn Logik und Hardware verwoben sind, kannst du nur am echten Gerät testen.
Jeder Test braucht Flashen, Verkabeln, manuelles Drücken -- langsam und fehleranfällig.

\textbf{Regel:} Logik muss isoliert testbar sein, Hardware-Pfad end-to-end.

\begin{table}[H]
\centering
\caption{Test-Environments in platformio.ini}
\label{tab:test-environments}
\begin{tabularx}{\textwidth}{@{}l l X@{}}
\toprule
\textbf{Environment} & \textbf{Typ} & \textbf{Was wird getestet} \\
\midrule
\code{seeed\_xiao\_esp32s3\_test} & Unit & Nur \code{logic/} (Debounce, Selection, Pipeline) \\
\code{seeed\_xiao\_esp32s3\_e2e}  & E2E  & \code{logic/} + \code{drivers/} + \code{hal/} mit echter Hardware \\
\bottomrule
\end{tabularx}
\end{table}

\paragraph{Unit-Test-Beispiel (Debounce)}

\begin{lstlisting}[style=arduino,caption={Unit-Test für Debounce},label={lst:debounce-test}]
void test_debounce_filters_noise() {
    DebounceState state = {0};

    // Simuliere Prellen: 0 -> 1 -> 0 -> 1 -> 1 -> 1
    TEST_ASSERT_FALSE(debounce_update(&state, 1, 3));  // Zaehler = 1
    TEST_ASSERT_FALSE(debounce_update(&state, 0, 3));  // Reset
    TEST_ASSERT_FALSE(debounce_update(&state, 1, 3));  // Zaehler = 1
    TEST_ASSERT_FALSE(debounce_update(&state, 1, 3));  // Zaehler = 2
    TEST_ASSERT_TRUE(debounce_update(&state, 1, 3));   // Zaehler = 3 -> Edge!
}
\end{lstlisting}

\begin{infobox}[Konsequenz]
Unit-Tests laufen in Sekunden, nicht Minuten. Du kannst Debounce-Schwellenwerte
experimentell optimieren, ohne jedes Mal zu flashen.
\end{infobox}

% -----------------------------------------------------------------------------
\subsection{Änderungen haben Pflichtfolgen}
\label{subsec:pflichtfolgen}

\textbf{Regel:} Bestimmte Änderungen erzwingen bestimmte Checks.

\begin{table}[H]
\centering
\caption{Pflicht-Checks nach Änderungstyp}
\label{tab:pflicht-checks}
\begin{tabularx}{\textwidth}{@{}l X@{}}
\toprule
\textbf{Änderung} & \textbf{Pflicht} \\
\midrule
Pins/SPI/Mapping            & \code{wiring.md} aktualisieren + E2E-Test \\
Debounce/Selection/Pipeline & Unit-Tests (\code{seeed\_xiao\_esp32s3\_test}) \\
Treiber/HAL                 & E2E-Test + Timing-Messung \\
Counts (10 $\to$ 100)       & Unit-Tests + E2E-Test + Timing-Prüfung \\
\bottomrule
\end{tabularx}
\end{table}

\paragraph{Testbefehle}

\begin{lstlisting}[style=shell,caption={Testbefehle für Unit- und E2E-Tests},label={lst:testbefehle}]
# Unit-Tests (Logik)
pio test -e seeed_xiao_esp32s3_test -v

# E2E-Hardware-Test
pio test -e seeed_xiao_esp32s3_e2e -f test_e2e_hw -v
\end{lstlisting}

% -----------------------------------------------------------------------------
\subsection{Minimal-Checkliste}
\label{subsec:checkliste}

Bevor du einen Change committest:

\begin{itemize}
  \item[$\square$] SSOT nur in \code{include/hw/*}
  \item[$\square$] Datenfluss bleibt: read $\to$ update $\to$ write
  \item[$\square$] One-Hot entsteht in genau einem Pfad (clear $\to$ set)
  \item[$\square$] Letztes Byte maskiert (\code{*\_LAST\_MASK})
  \item[$\square$] \code{logic/} ist hardware-unabhängig
  \item[$\square$] Unit-Tests grün (\code{seeed\_xiao\_esp32s3\_test})
  \item[$\square$] Bei HW-relevanten Changes: E2E grün (\code{test\_e2e\_hw})
\end{itemize}

% -----------------------------------------------------------------------------
\subsection{Siehe auch}
\label{subsec:arch-siehe-auch}

Weiterführende Dokumentation:

\begin{itemize}
  \item \code{firmware/docs/architecture.md} -- Schichtenmodell-Details
  \item \code{firmware/docs/overview.md} -- Bit-Mapping (MSB-first)
  \item \code{firmware/docs/wiring.md} -- Hardware-Verdrahtung
  \item \code{apps/bridge/docs/PROTOCOL.md} -- Serial/WebSocket-Protokoll
\end{itemize}
